\documentclass[10pt,letterpaper,twocolumn]{article}

\usepackage[utf8]{inputenc}
\usepackage[T1]{fontenc}
\usepackage{newtxtext,newtxmath}
\usepackage[scaled=0.92]{helvet}
\usepackage[top=0.76in,bottom=0.82in,left=0.72in,right=0.72in,columnsep=0.23in]{geometry}
\usepackage{microtype}
\usepackage{amsmath,mathtools}
\usepackage{graphicx}
\usepackage{booktabs}
\usepackage{longtable}
\usepackage{array}
\usepackage{enumitem}
\usepackage{titlesec}
\usepackage{xcolor}
\usepackage[hidelinks]{hyperref}
\usepackage{url}
\usepackage{fancyhdr}
\usepackage{setspace}

\setstretch{0.995}
\setlength{\parindent}{0.9em}
\setlength{\parskip}{0pt}
\setlength{\columnsep}{0.23in}
\setlist[itemize]{leftmargin=1.2em,itemsep=0.2em,topsep=0.35em}
\setlist[enumerate]{leftmargin=1.3em,itemsep=0.2em,topsep=0.35em}
\setlength{\abovedisplayskip}{5pt plus 2pt minus 2pt}
\setlength{\belowdisplayskip}{5pt plus 2pt minus 2pt}
\setlength{\abovedisplayshortskip}{4pt plus 2pt minus 2pt}
\setlength{\belowdisplayshortskip}{4pt plus 2pt minus 2pt}
\setlength{\textfloatsep}{8pt plus 2pt minus 2pt}
\setlength{\floatsep}{7pt plus 2pt minus 2pt}
\setlength{\intextsep}{7pt plus 2pt minus 2pt}
\setlength{\emergencystretch}{1.5em}
\raggedbottom

\definecolor{TitleRule}{HTML}{D9DEE7}
\definecolor{TitleNote}{HTML}{4B5563}
\definecolor{HeadingInk}{HTML}{1F2937}
\definecolor{HeadingRule}{HTML}{CBD5E1}

\renewcommand{\thesection}{\arabic{section}}
\renewcommand{\thesubsection}{\thesection.\arabic{subsection}}
\renewcommand{\thesubsubsection}{\thesubsection.\arabic{subsubsection}}
\setcounter{secnumdepth}{2}

\titleformat{\section}
  {\large\sffamily\bfseries\color{HeadingInk}}
  {\thesection}
  {0.65em}
  {}
\titleformat{name=\section,numberless}
  {\large\sffamily\bfseries\color{HeadingInk}}
  {}
  {0pt}
  {}
\titleformat{\subsection}
  {\normalsize\sffamily\bfseries\color{HeadingInk}}
  {\thesubsection}
  {0.55em}
  {}
\titleformat{\subsubsection}
  {\normalsize\sffamily\bfseries\color{HeadingInk}}
  {\thesubsubsection}
  {0.5em}
  {}
\titlespacing*{\section}{0pt}{0.95em}{0.35em}
\titlespacing*{\subsection}{0pt}{0.75em}{0.25em}
\titlespacing*{\subsubsection}{0pt}{0.55em}{0.2em}

\pagestyle{fancy}
\fancyhf{}
\fancyfoot[C]{\small\thepage}
\renewcommand{\headrulewidth}{0pt}


\providecommand{\tightlist}{%
  \setlength{\itemsep}{0pt}\setlength{\parskip}{0pt}}

\title{Geodesic-Provenance Coherence Bounds for Fourier-Domain Photon-Ring Inference}
\author{Jonathan R. Landers}
\date{}

\newcommand{\affiliationnote}{Independent researcher.\\
This note extends the structured Fourier-domain analysis in \\
\emph{A Fourier-Domain Analysis of BHEX for Photon-Ring Inference} and is not affiliated\\
with the BHEX collaboration or mission team.}

\makeatletter
\renewcommand{\maketitle}{%
  \twocolumn[{%
    \begin{@twocolumnfalse}
      \vspace*{-1.3em}
      {\LARGE\sffamily\bfseries\raggedright \@title \par}
      \vspace{0.35em}
      {\normalsize\sffamily\raggedright \@author \par}
      \vspace{0.16em}
      {\small\color{TitleNote}\raggedright \affiliationnote \par}
      \vspace{0.65em}
      {\color{TitleRule}\hrule height 0.8pt}
      \vspace{0.8em}
    \end{@twocolumnfalse}
  }]
}
\makeatother

\begin{document}
\maketitle

In the companion note \emph{A Fourier-Domain Analysis of BHEX for Photon-Ring Inference}
[\hyperlink{ref-1}{1}], the visibility data are modeled as
\[
y = \alpha g_\theta + q + \varepsilon,
\]
where $g_\theta$ is a photon-ring visibility family, $q$ is structured nuisance,
and $\varepsilon$ is noise. The main conclusion there is conceptual and useful: the
long-baseline amplitude heuristic can lose \textbf{readability} before the ring loses
\textbf{recoverability} under a structured inverse-problem formulation [\hyperlink{ref-1}{1}].

The next natural question is whether the nuisance term can be constrained more sharply by physics.
The purpose of this note is to show that the answer is yes: if one lifts the model from visibility
space back to photon initial-condition space and keeps track of \textbf{geodesic provenance},
then the relevant ring-background coherence admits an explicit upper bound in terms of the
cross-correlation of escaping geodesic families. This turns the sharpening from a structural remark
into a quantitative theorem, still in the spirit of the BHEX long-baseline program
[\hyperlink{ref-2}{2}, \hyperlink{ref-3}{3}, \hyperlink{ref-4}{4}].

\section{Original formalism and provenance lift}

The original abstraction is
\[
y = \alpha g_\theta + q + \varepsilon,
\]
with the nuisance treated as an abstract structured class. To refine this, let
\[
z = (p,\xi)
\]
denote a photon initial condition, where $p$ is an emission event and $\xi$ is an initial null direction.
Write $\gamma_z$ for the associated null geodesic and split phase space into
\[
E = \{z : \gamma_z \text{ escapes}\}, \qquad C = \{z : \gamma_z \text{ is captured}\}.
\]
Only $E$ contributes to the observed data. If $f$ is an emission density over initial conditions,
then the observed image can be written schematically as
\[
I = T_E f,
\]
where $T_E$ pushes escaping photons to the image plane, and the measured visibility is
\[
y = \mathcal F(T_E f) + \varepsilon.
\]

Now split the escaping set into two provenance classes,
\[
E = E_{\mathrm{ring}} \,\dot\cup\, E_{\mathrm{bg}},
\]
where $E_{\mathrm{ring}}$ consists of near-critical escaping geodesics that generate the photon-ring signal,
and $E_{\mathrm{bg}}$ contains the remaining escaping geodesics. Then
\[
T_E f = T_{E_{\mathrm{ring}}} f + T_{E_{\mathrm{bg}}} f.
\]
Defining
\[
h_\theta := T_{E_{\mathrm{ring}}} f_\theta, \qquad b := T_{E_{\mathrm{bg}}} f,
\]
gives
\[
I = \alpha h_\theta + b, \qquad y = \alpha g_\theta + q + \varepsilon,
\]
with
\[
g_\theta = \mathcal F h_\theta, \qquad q = \mathcal F b.
\]
The functional form of the inverse problem is unchanged, but the nuisance class is no longer arbitrary:
it is generated by ordinary escaping geodesics.

\section{Operator setup and provenance-constrained coherence}

Let $\mathcal Z$ be photon initial-condition space equipped with a measure $\nu$, and let $\mathcal H$
be the visibility Hilbert space. Define two bounded linear operators
\[
A_r := \mathcal F T_{E_{\mathrm{ring}}}, \qquad A_b := \mathcal F T_{E_{\mathrm{bg}}},
\]
so that ring-generated and background-generated visibilities are of the form $A_r f$ and $A_b h$,
respectively.

Fix an admissible class $\mathcal A$ of background emission densities and let the provenance-constrained
nuisance family be
\[
Q_{\mathrm{prov}} := \{A_b h : h \in \mathcal A,\ A_b h \neq 0\}.
\]
For a ring template $g \in \mathcal H$, define the family coherence
\[
\mu(g,Q_{\mathrm{prov}})
:=
\sup_{q\in Q_{\mathrm{prov}}}
\frac{|\langle g,q\rangle|}{\|g\|\,\|q\|}.
\]
When $g = A_r f_\theta$, this is the physically relevant ring-background overlap.

\par\smallskip\noindent{}\textbf{Theorem (Provenance-based coherence bound).}
Assume there exist constants $\sigma_r,\sigma_b>0$ such that
\[
\|A_r f\| \ge \sigma_r \|f\|, \qquad \|A_b h\| \ge \sigma_b \|h\|
\]
for all admissible ring coefficients $f$ and background coefficients $h \in \mathcal A$.
Let $g = A_r f_\theta$. Then
\[
\mu(g,Q_{\mathrm{prov}})
\le
\frac{\|A_b^*A_r\|}{\sigma_r\sigma_b}.
\]
Moreover, if $A_r$ and $A_b$ admit a visibility kernel $K(u,z)$ and one defines the cross-Gram kernel
\[
G(z,z') := \langle K(\cdot,z),K(\cdot,z')\rangle_{\mathcal H},
\]
then
\begin{equation}
\begin{aligned}
\mu(g,Q_{\mathrm{prov}})
&\le
\frac{1}{\sigma_r\sigma_b}
\biggl(
\int_{E_{\mathrm{ring}}}\!\int_{E_{\mathrm{bg}}}
|G(z,z')|^2\,d\nu(z)\,d\nu(z')
\biggr)^{1/2}.
\end{aligned}
\end{equation}

\par\smallskip\noindent{}\textbf{Proof.} Let $q=A_b h \in Q_{\mathrm{prov}}$. Then
\[
\begin{aligned}
|\langle g,q\rangle|
&= |\langle A_r f_\theta, A_b h\rangle| \\
&= |\langle A_b^*A_r f_\theta, h\rangle| \\
&\le \|A_b^*A_r\|\,\|f_\theta\|\,\|h\|.
\end{aligned}
\]
By the lower bounds,
\[
\|g\|=\|A_r f_\theta\|\ge \sigma_r\|f_\theta\|,
\qquad
\|q\|=\|A_b h\|\ge \sigma_b\|h\|.
\]
Hence
\[
\frac{|\langle g,q\rangle|}{\|g\|\,\|q\|}
\le
\frac{\|A_b^*A_r\|}{\sigma_r\sigma_b}.
\]
Taking the supremum over $q \in Q_{\mathrm{prov}}$ proves the first claim. The second follows from the Hilbert-Schmidt bound
\[
\begin{aligned}
\|A_b^*A_r\|
&\le \|A_b^*A_r\|_{\mathrm{HS}} \\
&= \left(
\int_{E_{\mathrm{ring}}}\!\int_{E_{\mathrm{bg}}}
|G(z,z')|^2\,d\nu(z)\,d\nu(z')
\right)^{1/2}.
\end{aligned}
\]
\hfill$\square$

This is the new mathematical content. The coherence term that controlled the mismatch bound in the
original note [\hyperlink{ref-1}{1}] is now itself controlled by a cross-correlation quantity on photon
initial-condition space.

\section{A geometric separation corollary}

The preceding theorem becomes more informative if the two escaping families are separated by a geometric
criticality index. Let $\chi(z)$ be any scalar measure of ring-likeness along an escaping geodesic,
for example a winding count surrogate, a residence-time proxy near the photon region, or another criticality
parameter.

Assume there exists a threshold $\tau$ and a gap $\Delta>0$ such that
\[
\chi(z) \ge \tau \quad \text{for } z\in E_{\mathrm{ring}},
\qquad
\chi(z') \le \tau-\Delta \quad \text{for } z'\in E_{\mathrm{bg}}.
\]
Assume also that the cross-Gram kernel decays with criticality separation in the sense that
\[
|G(z,z')| \le M e^{-\beta |\chi(z)-\chi(z')|}
\]
for some $M,\beta>0$.

\par\smallskip\noindent{}\textbf{Corollary (Exponential coherence decay under geodesic separation).}
Under the preceding assumptions,
\begin{equation}
\begin{aligned}
\mu(g,Q_{\mathrm{prov}})
&\le
\frac{M}{\sigma_r\sigma_b} e^{-\beta\Delta} \\
&\quad \cdot \sqrt{\nu(E_{\mathrm{ring}})\,\nu(E_{\mathrm{bg}})}.
\end{aligned}
\end{equation}

\par\smallskip\noindent{}\textbf{Proof.} The separation assumption implies
\[
|\chi(z)-\chi(z')| \ge \Delta
\quad
\text{for } z\in E_{\mathrm{ring}},\ z'\in E_{\mathrm{bg}},
\]
so the kernel bound gives
\[
|G(z,z')| \le M e^{-\beta\Delta}.
\]
Substituting into the Hilbert-Schmidt estimate of the theorem yields
\begin{equation}
\begin{aligned}
\biggl(
\int_{E_{\mathrm{ring}}}\!\int_{E_{\mathrm{bg}}}
|G(z,z')|^2\,d\nu(z)\,d\nu(z')
\biggr)^{1/2}
&\le M e^{-\beta\Delta} \\
&\quad \cdot \sqrt{\nu(E_{\mathrm{ring}})\,\nu(E_{\mathrm{bg}})},
\end{aligned}
\end{equation}
which proves the claim.
\hfill$\square$

This is a stronger statement than monotonicity under subset restriction. It says that if ring-generating
escape and ordinary escape are geometrically separated, then the Fourier-domain coherence decays
\emph{quantitatively}. The sharper bound comes from phase-space geometry, not merely from saying that the
nuisance class is smaller.

\section{Sharpening the original mismatch and recoverability story}

In the companion note [\hyperlink{ref-1}{1}], the heuristic-versus-structured mismatch is controlled by a term of the form
\[
\mu(g,\mathcal P)\frac{\|p\|}{\|g\|}
\]
plus noise. The new theorem suggests replacing an abstract nuisance overlap by an explicit provenance bound:
\begin{equation}
\mu(g,Q_{\mathrm{prov}})\frac{\|q\|}{\|g\|}
\le
\frac{\|A_b^*A_r\|}{\sigma_r\sigma_b}\frac{\|q\|}{\|g\|}.
\end{equation}
and, under geodesic separation,
\begin{equation}
\begin{aligned}
\mu(g,Q_{\mathrm{prov}})\frac{\|q\|}{\|g\|}
&\le
\frac{M}{\sigma_r\sigma_b}e^{-\beta\Delta} \\
&\quad \cdot \sqrt{\nu(E_{\mathrm{ring}})\,\nu(E_{\mathrm{bg}})}\frac{\|q\|}{\|g\|}.
\end{aligned}
\end{equation}
So the old mismatch story becomes sharper in a specific way: the bias term is small whenever
near-critical escaping flow and ordinary escaping flow are sufficiently decorrelated.

The recoverability proposition from the original note remains conceptually the same, but its nuisance class is
now physically constrained. The point is not simply that captured photons do not appear in the data. The more
important fact is that the \emph{observable} nuisance must come from ordinary escaping geodesics, and that
restriction can now be turned into a quantitative visibility-space overlap bound.

\section{Physical interpretation and relation to BHEX}

The BHEX idea is to use sufficiently long baselines to isolate the universal photon-ring signal from more
complicated astrophysical emission [\hyperlink{ref-2}{2}, \hyperlink{ref-3}{3}]. The companion note [\hyperlink{ref-1}{1}] recasts that idea as a Fourier-domain structured source-separation problem. The present note pushes one step further: it identifies a route by which the
relevant ring-background coherence might actually be estimated from geodesic geometry.

In practical terms, the theorem says that the Fourier-domain nuisance is controlled by how strongly the
ring-generating escaping family and the background escaping family correlate after transport to visibility
space. If those families are dynamically distinct, then structured separation has a principled advantage beyond
a vague “smaller nuisance” story.

This remains an abstract framework. To turn it into an astrophysical forecast, one would need to specify
admissible emission classes, choose a concrete criticality index, and estimate the cross-Gram kernel from ray
tracing or simulation. But that is exactly the kind of bridge from abstract mathematics to physically informed
BHEX-style inference that the original note pointed toward [\hyperlink{ref-1}{1}, \hyperlink{ref-3}{3}, \hyperlink{ref-4}{4}].

\section{Key interesting points}

\begin{itemize}
\tightlist
\item This note keeps the same Fourier-domain inverse-problem skeleton as the companion note, but lifts the
nuisance model back to photon initial-condition space.
\item The key new theorem bounds the ring-background coherence by $\|A_b^*A_r\|/(\sigma_r\sigma_b)$, and then by a
cross-Gram integral over escaping geodesic families.
\item Under a geometric separation hypothesis in a criticality index, the coherence decays exponentially in the
separation gap $\Delta$.
\item This is stronger than saying the nuisance class is merely smaller; it gives a concrete mechanism by which
provenance can sharpen the mismatch term in the original note.
\item The new note therefore earns independence only to the extent that coherence itself is now bounded by phase-space geometry rather than treated as an abstract parameter.
\end{itemize}

\section*{References}
\addcontentsline{toc}{section}{References}

\noindent{}\hypertarget{ref-1}{}{[}1{]} Jonathan R. Landers,
\emph{A Fourier-Domain Analysis of BHEX for Photon-Ring Inference}, 2026.

\noindent{}\hypertarget{ref-2}{}{[}2{]} Black Hole Explorer Collaboration,
\emph{Black Hole Explorer: Motivation and Vision}, arXiv:2406.12917.
\url{https://arxiv.org/abs/2406.12917}

\noindent{}\hypertarget{ref-3}{}{[}3{]} Black Hole Explorer Collaboration,
\emph{The Black Hole Explorer: Photon Ring Science, Detection and Shape Measurement},
arXiv:2406.09498. \url{https://arxiv.org/abs/2406.09498}

\noindent{}\hypertarget{ref-4}{}{[}4{]} Black Hole Explorer Collaboration,
\emph{Interferometric Inference of Black Hole Spin from Photon Ring Size and Brightness},
arXiv:2509.23628. \url{https://arxiv.org/abs/2509.23628}

\end{document}
